% GENERATED FILE. Edit canonical lessons, then run npm run generate:books.
% canonical-source-hash: fnv1a64:4221ec45e549345f
% canonical-lessons: SA-C334-raksasah, SA-C334-asurah, SA-C334-vadhah, SA-C334-hanti, SA-C334-nigirati

\chapter{Gods and Demons}
\label{ch:sa-gods-and-demons}
\hlchaptermodality{334}

\begin{chapteropening}
\textbf{By the end of this chapter:} \emph{I can tell a story of gods and demons.}

Complete the last lesson of chapter 334: I can tell a story of gods and demons.
\end{chapteropening}

\section[\texorpdfstring{\sk{राक्षसः} (rākṣasaḥ)}{राक्षसः (rākṣasaḥ)}]{\sk{राक्षसः} (rākṣasaḥ) — a demon}

\label{lesson:SA-C334-raksasah}



\begin{quote}
Before the new one: say the Sanskrit for shows favour, then the Sanskrit for immortal.
\end{quote}

\subsection*{You'll want to know: \sk{राक्षसः}}
\textbf{\sk{राक्षसः}} — \emph{rākṣasaḥ} — \textquotedblleft{}a demon\textquotedblright{}.

\begin{quote}
\sk{वने} \sk{एकः} \sk{भयंकरः} \sk{राक्षसः} \sk{वसति।}
\end{quote}

\begin{quote}
\emph{vane ekaḥ bhayaṅkaraḥ rākṣasaḥ vasati.}
\end{quote}

\begin{quote}
A terrible demon lives in the forest.
\end{quote}

A masculine noun.

\begin{cousinweb}
\textbf{\sk{राक्षसः}} (\emph{rākṣasaḥ}) is built on \emph{rakṣas}, \textquotedblleft{}a demon, an evil being\textquotedblright{}; MW gives \emph{rākṣasa} as belonging to the Rakṣas and, as a noun, a demon (MW).
\end{cousinweb}

\begin{grammarlens}[title={What you've built}]
One more word for telling what happened.
\end{grammarlens}

\subsection*{Your turn}
\begin{itemize}
\raggedright
  \item \emph{Say it:} \emph{rākṣasaḥ}
  \item \emph{Say it:} \emph{rākṣasaḥ}, once more
  \item \emph{Say it:} \emph{vane ekaḥ bhayaṅkaraḥ rākṣasaḥ vasati.} — A terrible demon lives in the forest.
  \item \emph{Recall:} say the Sanskrit for shows favour, then the Sanskrit for immortal, then say \emph{rākṣasaḥ} again
\end{itemize}

\begin{tcolorbox}[breakable,colback=teal!4,colframe=teal!35!black,title={Before you move on}]
What does \sk{राक्षसः} mean? (\textquotedblleft{}A demon\textquotedblright{}.) Say it once more.
\end{tcolorbox}
\section[\texorpdfstring{\sk{असुरः} (asuraḥ)}{असुरः (asuraḥ)}]{\sk{असुरः} (asuraḥ) — an asura; an enemy of the gods}

\label{lesson:SA-C334-asurah}



\begin{quote}
Before the new one: say the Sanskrit for immortal, then the Sanskrit for a demon.
\end{quote}

\subsection*{You'll want to know: \sk{असुरः}}
\textbf{\sk{असुरः}} — \emph{asuraḥ} — \textquotedblleft{}an asura; an enemy of the gods\textquotedblright{}.

\begin{quote}
\sk{देवाः} \sk{असुरान्} \sk{जयन्ति।}
\end{quote}

\begin{quote}
\emph{devāḥ asurān jayanti.}
\end{quote}

\begin{quote}
The gods defeat the asuras.
\end{quote}

A masculine noun.

\begin{cousinweb}
\textbf{\sk{असुरः}} (\emph{asuraḥ}) MW reports a grammarians' derivation from √as and, beside older senses of a spirit, gives \emph{asura} as an evil spirit, an opponent of the gods, so its origin is uncertain (MW).
\end{cousinweb}

\begin{grammarlens}[title={What you've built}]
One more word for telling what happened.
\end{grammarlens}

\subsection*{Your turn}
\begin{itemize}
\raggedright
  \item \emph{Say it:} \emph{asuraḥ}
  \item \emph{Say it:} \emph{asuraḥ}, once more
  \item \emph{Say it:} \emph{devāḥ asurān jayanti.} — The gods defeat the asuras.
  \item \emph{Recall:} say the Sanskrit for immortal, then the Sanskrit for a demon, then say \emph{asuraḥ} again
\end{itemize}

\begin{tcolorbox}[breakable,colback=teal!4,colframe=teal!35!black,title={Before you move on}]
What does \sk{असुरः} mean? (\textquotedblleft{}An asura; an enemy of the gods\textquotedblright{}.) Say it once more.
\end{tcolorbox}
\section[\texorpdfstring{\sk{वधः} (vadhaḥ)}{वधः (vadhaḥ)}]{\sk{वधः} (vadhaḥ) — the slaying; a killing}

\label{lesson:SA-C334-vadhah}



\begin{quote}
Before the new one: say the Sanskrit for a demon, then the Sanskrit for an enemy of the gods.
\end{quote}

\subsection*{You'll want to know: \sk{वधः}}
\textbf{\sk{वधः}} — \emph{vadhaḥ} — \textquotedblleft{}the slaying; a killing\textquotedblright{}.

\begin{quote}
\sk{वीरः} \sk{राक्षसस्य} \sk{वधं} \sk{करोति।}
\end{quote}

\begin{quote}
\emph{vīraḥ rākṣasasya vadhaṃ karoti.}
\end{quote}

\begin{quote}
The hero slays the demon.
\end{quote}

A masculine noun.

\begin{cousinweb}
\textbf{\sk{वधः}} (\emph{vadhaḥ}) MW gives \emph{vadha} as one who kills, a deadly weapon, and the act of striking or killing, without a root in the entry, so its origin is left uncertain here (MW).
\end{cousinweb}

\begin{grammarlens}[title={What you've built}]
One more word for telling what happened.
\end{grammarlens}

\subsection*{Your turn}
\begin{itemize}
\raggedright
  \item \emph{Say it:} \emph{vadhaḥ}
  \item \emph{Say it:} \emph{vadhaḥ}, once more
  \item \emph{Say it:} \emph{vīraḥ rākṣasasya vadhaṃ karoti.} — The hero slays the demon.
  \item \emph{Recall:} say the Sanskrit for a demon, then the Sanskrit for an enemy of the gods, then say \emph{vadhaḥ} again
\end{itemize}

\begin{tcolorbox}[breakable,colback=teal!4,colframe=teal!35!black,title={Before you move on}]
What does \sk{वधः} mean? (\textquotedblleft{}The slaying; a killing\textquotedblright{}.) Say it once more.
\end{tcolorbox}
\section[\texorpdfstring{\sk{हन्ति} (hanti)}{हन्ति (hanti)}]{\sk{हन्ति} (hanti) — kills; slays}

\label{lesson:SA-C334-hanti}



\begin{quote}
Before the new one: say the Sanskrit for an enemy of the gods, then the Sanskrit for the slaying.
\end{quote}

\subsection*{You'll want to know: \sk{हन्ति}}
\textbf{\sk{हन्ति}} — \emph{hanti} — \textquotedblleft{}kills; slays\textquotedblright{}.

\begin{quote}
\sk{वीरः} \sk{खड्गेन} \sk{राक्षसं} \sk{हन्ति।}
\end{quote}

\begin{quote}
\emph{vīraḥ khaḍgena rākṣasaṃ hanti.}
\end{quote}

\begin{quote}
The hero kills the demon with a sword.
\end{quote}

Past (he …): \sk{सः} \sk{राक्षसं} \sk{हतवान्}, \emph{saḥ rākṣasaṃ hatavān}.

\begin{cousinweb}
\textbf{\sk{हन्ति}} (\emph{hanti}) comes from the root √han, \textquotedblleft{}to strike, kill\textquotedblright{}; MW compares Greek \emph{theinō} and \emph{phonos} (MW).
\end{cousinweb}

\begin{grammarlens}[title={What you've built}]
One more word for telling what happened.
\end{grammarlens}

\subsection*{Your turn}
\begin{itemize}
\raggedright
  \item \emph{Say it:} \emph{hanti}
  \item \emph{Say it:} \emph{hanti}, once more
  \item \emph{Say it:} \emph{vīraḥ khaḍgena rākṣasaṃ hanti.} — The hero kills the demon with a sword.
  \item \emph{Recall:} say the Sanskrit for an enemy of the gods, then the Sanskrit for the slaying, then say \emph{hanti} again
\end{itemize}

\begin{tcolorbox}[breakable,colback=teal!4,colframe=teal!35!black,title={Before you move on}]
What does \sk{हन्ति} mean? (\textquotedblleft{}Kills; slays\textquotedblright{}.) Say it once more.
\end{tcolorbox}
\section[\texorpdfstring{\sk{निगिरति} (nigirati)}{निगिरति (nigirati)}]{\sk{निगिरति} (nigirati) — swallows}

\label{lesson:SA-C334-nigirati}



\begin{quote}
Before the new one: say the Sanskrit for the slaying, then the Sanskrit for kills.
\end{quote}

\subsection*{You'll want to know: \sk{निगिरति}}
\textbf{\sk{निगिरति}} — \emph{nigirati} — \textquotedblleft{}swallows\textquotedblright{}.

\begin{quote}
\sk{सर्पः} \sk{मूषकं} \sk{निगिरति।}
\end{quote}

\begin{quote}
\emph{sarpaḥ mūṣakaṃ nigirati.}
\end{quote}

\begin{quote}
The snake swallows the mouse.
\end{quote}

Past (it …): \sk{सर्पः} \sk{मूषकं} \sk{निगीर्णवान्}, \emph{sarpaḥ mūṣakaṃ nigīrṇavān}.

\begin{cousinweb}
\textbf{\sk{निगिरति}} (\emph{nigirati}) is built from \emph{ni}, \textquotedblleft{}down\textquotedblright{}, and the root √gṝ, \textquotedblleft{}to swallow\textquotedblright{}; MW gives \emph{ni-gṝ} as to swallow, gulp down (MW).
\end{cousinweb}

\begin{grammarlens}[title={What you've built}]
One more word for telling what happened.
\end{grammarlens}

\subsection*{Your turn}
\begin{itemize}
\raggedright
  \item \emph{Say it:} \emph{nigirati}
  \item \emph{Say it:} \emph{nigirati}, once more
  \item \emph{Say it:} \emph{sarpaḥ mūṣakaṃ nigirati.} — The snake swallows the mouse.
  \item \emph{Recall:} say the Sanskrit for the slaying, then the Sanskrit for kills, then say \emph{nigirati} again
\end{itemize}

\begin{tcolorbox}[breakable,colback=teal!4,colframe=teal!35!black,title={Before you move on}]
What does \sk{निगिरति} mean? (\textquotedblleft{}Swallows\textquotedblright{}.) Say it once more.
\end{tcolorbox}
